\chapter{Meta-tasks: Workflows of Other Engines}
\label{app:meta-tasks}

Lecture~22 composed several DAGonStar workflows with \texttt{DAG\_TPS}. Many
useful pipelines, however, are already written for another engine: Nextflow,
Snakemake, Parsl, StreamFlow or a CWL runner. DAGonStar~3.0 runs such a
workflow, unmodified, as one task of type \texttt{TaskType.META}. Data moves in
and out through ordinary \texttt{workflow://} references, so the external
workflow is a node of the DAGonStar graph like any other task. This appendix
accompanies course deck~A3; its sources are \texttt{dagon/meta/} and
\texttt{docs/meta\_tasks.md} of the DAGonStar repository.

\section{Two Properties by Construction}

\begin{definitionbox}{Meta-task}
A task that runs a complete workflow written for another engine through the
engine's documented command line. The engine is never modified: DAGonStar
writes the engine's own parameter files (a Nextflow \texttt{-params-file}, a
Snakemake \texttt{--configfile}, a CWL job document, a Python call for Parsl).
Every engine-specific detail lives in one \emph{adapter}; staging, output
collection, checkpoints, FAIR records, JSON and CWL behave the same for all
engines.
\end{definitionbox}

\begin{table}[h]
\centering
\begin{tabular}{@{}lll@{}}
\toprule
\texttt{engine} & What runs & Inputs reach the workflow as \\
\midrule
\texttt{nextflow} & \texttt{nextflow run} with a params file & \texttt{params.<name>} \\
\texttt{snakemake} & \texttt{snakemake --configfile} & \texttt{config["<name>"]} \\
\texttt{parsl} & a Python function or script & keyword arguments \\
\texttt{streamflow} & \texttt{streamflow run} & a CWL job document \\
\texttt{cwl} & \texttt{cwltool} or another CWL runner & a CWL job document \\
\texttt{command} & any argument vector & a JSON parameter file \\
\texttt{mock} & no process & placeholders (portable emulation) \\
\bottomrule
\end{tabular}
\caption{The built-in meta-task adapters of DAGonStar~3.0.0.}
\label{tab:meta-adapters}
\end{table}

\section{The Task}

A meta-task takes an \texttt{engine}, a \texttt{workflow} (a local file or
directory, or, for engines that resolve names themselves, a name such as
\texttt{nextflow-io/hello}), \texttt{inputs}, \texttt{outputs} and
\texttt{engine\_options}. An input whose value is a \texttt{workflow://} string
is a staged file or directory; any other JSON value is passed through. An
output maps a name to a path below \texttt{outputs/}, or to a descriptor with
\texttt{path}, \texttt{source} (where the engine writes, relative to its
working directory \texttt{run/}) and \texttt{engine\_output} (an output the
engine reports by name). Options the adapter does not declare are rejected, and
so are keys that look like credentials: every validation error is a
\texttt{MetaValidationError}, raised when the task is created.

The lifecycle has five steps. A dependency is added for each producer named by
an input; inputs are staged to \texttt{inputs/<name>/}; the adapter writes the
engine's parameter files under \texttt{.dagon/} and runs the engine in
\texttt{run/}, without a shell and with the task timeout; each declared output
is copied from \texttt{run/<source>} to \texttt{outputs/<path>} and hashed with
SHA-256; and \texttt{.dagon/meta\_result.json} records the engine, the workflow
identity, the digests and the duration. A failure is classified in
\texttt{task.failure\_phase} as \texttt{validation}, \texttt{engine},
\texttt{staging}, \texttt{execution} or \texttt{outputs}, and a failed meta-task
fails its dependants.

\section{Three Engines on the Station Data}

The course example has two real engines and one emulated. The first is a
stand-in program that knows nothing about DAGonStar (\Cref{lst:station-engine}),
driven by the generic \texttt{command} adapter. The second is a CWL tool
(\Cref{lst:count-records}), run by \texttt{cwltool} when it is installed. The
third is a Nextflow pipeline under portable emulation.

\lstinputlisting[caption={A stand-in engine (\texttt{examples/A3-meta/engines/station\_engine.py}).},label={lst:station-engine}]{../../examples/A3-meta/engines/station_engine.py}

\lstinputlisting[style=bash,caption={A CWL tool (\texttt{examples/A3-meta/engines/count\_records.cwl}).},label={lst:count-records}]{../../examples/A3-meta/engines/count_records.cwl}

\lstinputlisting[caption={Meta-tasks on the station data (\texttt{examples/A3-meta/meta\_task.py}).},label={lst:meta-task}]{../../examples/A3-meta/meta_task.py}

\lstinputlisting[style=bash,caption={Recorded output of \texttt{meta\_task.py}.}]{../../examples/A3-meta/output/meta_task.txt}

The means match those of Appendix~\ref{app:volumes} and Lecture~12, and the CWL
tool counts the $3 \times 4 = 12$ data records of the dataset.

\section{Checkpoints and Evidence}

A completed meta-task is reused only when everything can be verified again:
the effective specification has the same SHA-256; the workflow definition is
local and has the same digest (a remote workflow always runs again); every
input, re-read from its producer's current folder, has the recorded digest;
and every collected output still exists with its recorded digest. A successful
folder alone is never enough, and checkpoints of emulated and real runs are
never mixed. The engine's own resume mechanism, such as Nextflow's
\texttt{resume}, is independent and opt-in through \texttt{engine\_options}.

\begin{practicebox}
Site defaults go into \texttt{dagon.ini}, in one section per engine:
\texttt{[meta.nextflow]} with \texttt{executable}, \texttt{[meta.parsl]} with
\texttt{python}. Task options take precedence, and the effective values are part
of the checkpoint identity.
\end{practicebox}

Under portable emulation the \texttt{mock} adapter starts no engine and writes
one JSON placeholder per declared output, marked as emulated in
\texttt{meta\_result.json}, in the checkpoint and in FAIR records. With FAIR
recording on, each collected output is an artifact with its checksum, and the
RO-Crate action of the task names the external workflow as its instrument.
\texttt{saveAsCWL} exports a meta-task as a \texttt{CommandLineTool} that runs
\texttt{dagon-meta-runner}, with a \texttt{dagon:MetaTask} hint; the engine and
the workflow definition must then be installed on the CWL execution host.

\begin{cautionbox}[What a meta-task does not see]
DAGonStar records the inputs, outputs and identity of the external workflow,
not what happens inside it. The engine's own provenance, such as a Nextflow
report or the StreamFlow database, stays in the task's scratch folder.
\end{cautionbox}
